\section{Conclusions}\label{sec:conclusion}
We provide the first nationwide evidence on how transitioning from decentralized to centralized school assignment shapes students' subsequent behavior and outcomes in a separate centralized college admissions process. Leveraging the staggered and partially randomized rollout of Chile's \textit{Sistema de Admisión Escolar} (SAE), we estimate dynamic treatment effects that isolate the causal impact of centralized school choice on the full pipeline from secondary schooling to higher education. Two findings stand out. First, exposure to centralized school assignment substantially increases the likelihood that students apply to, are admitted to, and enroll in programs within the centralized college admissions system. Second, these gains are concentrated among students historically underserved by decentralized admissions---female students, students from socioeconomically vulnerable schools, and students from lower-performing schools---indicating that centralized assignment can meaningfully reduce downstream inequities in access to higher education.

Beyond documenting average effects, our analysis sheds light on the mechanisms through which centralized school choice affects postsecondary trajectories. We show that the improvements are driven by increased admissions to less selective and historically under-subscribed programs, with no evidence of displacement in highly selective programs. This pattern suggests that centralized assignment expands participation without generating congestion in the most competitive segments of the market. The absence of effects on high-school GPA and entrance-exam performance further indicates that the reform operates primarily through changes in information, expectations, and application behavior rather than through improvements in academic preparation. Together, these results highlight the importance of institutional design in shaping how students navigate sequential matching environments and demonstrate that early exposure to transparent, standardized assignment processes can influence later decisions in related markets.

From a policy perspective, our findings underscore the broader value of coordinated choice systems in education markets. Centralized assignment not only reduces inefficiencies and inequities in school admissions but also appears to alter students' engagement with subsequent centralized processes, expanding access to higher education for groups that have historically faced barriers in decentralized settings. These insights are particularly relevant for policymakers considering reforms to school or college admissions systems, as they suggest that the benefits of centralized assignment extend beyond the immediate allocation problem and can propagate across educational stages. More broadly, our results point to promising avenues for future work on how early institutional experiences shape behavior in later matching environments, how information and transparency affect students' strategic decisions, and how coordinated systems can be designed to promote equity and efficiency across the educational pipeline.


